\chapter{El principio de contracción y sus aplicaciones}\label{chap:11}
\section{Resolver una ecuación mediante iteraciones}
Una ecuación puede escribirse en la forma $x=T(x)$. Si aplicar $T$ reduce las distancias por un factor fijo menor que uno, repetir la operación produce aproximaciones cada vez más estables. La completitud garantiza que esas aproximaciones desembocan en un punto del espacio. Además de existencia y unicidad, el argumento proporciona cotas explícitas del error.

\begin{definicion}
Una \emph{contracción} de un espacio métrico $X$ es una aplicación $T:X\to X$ para la que existe $0\leq q<1$ con
\[
 d(Tx,Ty)\leq qd(x,y)\qquad(x,y\in X).
\]
Un \emph{punto fijo} es un punto $x_*$ que satisface $T(x_*)=x_*$.
\end{definicion}
\index{contraccion@contracción}\index{punto fijo!teorema de Banach}
El caso $q=0$ incluye las aplicaciones constantes. La desigualdad estricta $d(Tx,Ty)<d(x,y)$ para $x\ne y$ no proporciona necesariamente una constante uniforme $q<1$; el teorema siguiente exige dicha constante.

\begin{teorema}[Punto fijo de Banach]\label{ch11:banach}
Si $X$ es completo y no vacío y $T:X\to X$ es una contracción con constante $q<1$, entonces existe un único punto fijo $x_*$. Para cualquier $x_0\in X$, las iteraciones $x_{n+1}=T(x_n)$ convergen a él. Si $0<q<1$, se tienen las cotas
\[
 d(x_n,x_*)\leq\frac{q^n}{1-q}d(x_1,x_0),
 \qquad d(x_n,x_*)\leq\frac{q}{1-q}d(x_n,x_{n-1})\quad(n\geq1).
\]
\end{teorema}
\begin{proof}
Si $q=0$, $T$ es constante, digamos $T(x)=c$, y $T(c)=c$; además ningún otro punto puede ser fijo. Supongamos $0<q<1$. Por inducción,
\[
 d(x_{n+1},x_n)\leq q^nd(x_1,x_0).
\]
Para $m>n$, sumamos distancias entre términos consecutivos:
\[
 d(x_m,x_n)\leq\sum_{k=n}^{m-1}d(x_{k+1},x_k)
 \leq d(x_1,x_0)\sum_{k=n}^{m-1}q^k
 \leq\frac{q^n}{1-q}d(x_1,x_0).
\]
El extremo derecho tiende a cero y no depende de $m$, así que $(x_n)$ es de Cauchy. Sea $x_*$ su límite en $X$. Una contracción es Lipschitz y continua; luego $T(x_n)\to T(x_*)$. Pero $T(x_n)=x_{n+1}\to x_*$, y la unicidad del límite da $T(x_*)=x_*$. Si $y_*$ fuese otro punto fijo,
\[
 d(x_*,y_*)=d(Tx_*,Ty_*)\leq qd(x_*,y_*),
\]
de donde $(1-q)d(x_*,y_*)\leq0$. Como $1-q>0$, los puntos coinciden.

Pasando $m\to\infty$ en la estimación de las colas obtenemos la primera cota. Para la segunda repetimos la suma desde $x_n$ y usamos $d(x_{n+j+1},x_{n+j})\leq q^{j+1}d(x_n,x_{n-1})$; la serie geométrica suma $q/(1-q)$. Esta última cota permite estimar el error a partir de dos iteraciones ya calculadas.
\end{proof}

\begin{ejemplo}[Existencia por continuidad y existencia por contracción]
Toda función continua $f:[0,1]\to[0,1]$ tiene al menos un punto fijo. En efecto, $h(x)=x-f(x)$ satisface $h(0)\leq0$ y $h(1)\geq0$; el teorema del valor intermedio proporciona un cero. La unicidad no se deduce: la identidad tiene todos los puntos fijos. Por otro lado, si $f:\R\to\R$ es derivable y $|f'(x)|\leq q<1$ para todo $x$, el valor medio da $|f(x)-f(y)|\leq q|x-y|$. Banach produce entonces un único punto fijo y convergencia desde cualquier valor inicial.
\end{ejemplo}

\section{Sistemas no lineales y elección de una norma}
\begin{proposicion}\label{ch11:jacobiano}
Sea $R\subset\R^n$ un rectángulo cerrado no vacío y $T:R\to R$ una aplicación de clase $C^1$ en una vecindad de $R$. Si
\[
 \sup_{x\in R}\max_i\sum_{j=1}^n
 \left|\frac{\partial T_i}{\partial x_j}(x)\right|\leq q<1,
\]
entonces $T$ es una contracción para la norma máxima y tiene un único punto fijo en $R$.
\end{proposicion}
\begin{proof}
El segmento entre $x,y\in R$ permanece en $R$ por convexidad. Aplicando el teorema fundamental del cálculo a $t\mapsto T_i(y+t(x-y))$,
\[
 T_i(x)-T_i(y)=\int_0^1\sum_j
 \frac{\partial T_i}{\partial x_j}(y+t(x-y))(x_j-y_j)\,dt.
\]
El valor absoluto está acotado por $q\norm{x-y}_\infty$. Tomar el máximo en $i$ demuestra la contracción. El rectángulo es cerrado en el espacio completo $\R^n$, por lo que es completo; ahora se aplica el teorema~\ref{ch11:banach}. Además de las derivadas, debe comprobarse que $T(R)\subset R$. Sin esta condición las iteraciones podrían salir del dominio.
\end{proof}
En particular, basta que cada derivada parcial tenga valor absoluto no mayor que $q/n$. La hipótesis de rectángulo cerrado corrige la posible pérdida de completitud en un rectángulo abierto.

\begin{ejemplo}[Un sistema de tres ecuaciones]\label{ch11:sistema3}
Consideremos
\begin{align*}
 3x-\cos(yz)-\tfrac12&=0,\\
 x^2-81(y+0.1)^2+\sin z+1.06&=0,\\
 e^{-xy}+20z+\tfrac{10\pi-3}{3}&=0.
\end{align*}
La rama $y+0.1\geq0$ conduce a
\[
 T(x,y,z)=\left(\frac{\cos(yz)}3+\frac16,
 \frac{\sqrt{x^2+\sin z+1.06}}9-0.1,
 -\frac{e^{-xy}}{20}-\frac{10\pi-3}{60}\right).
\]
Verificamos todas las condiciones en $R=[-1,1]^3$. El radicando $Q=x^2+\sin z+1.06$ satisface
\[
 Q\geq1.06-\sin1>0.21,\qquad Q<2.91<4.
\]
La raíz está bien definida y es diferenciable en una vecindad del cubo. La primera coordenada de $T$ está entre $-1/6$ y $1/2$; la segunda está entre $-0.1$ y $2/9-0.1$, y la tercera satisface
\[
 |T_3|\leq e/20+(10\pi-3)/60<1.
\]
Así $T(R)\subset R$.

Las derivadas no nulas son
\begin{align*}
 \partial_yT_1&=-z\sin(yz)/3,&\partial_zT_1&=-y\sin(yz)/3,\\
 \partial_xT_2&=x/(9\sqrt Q),&\partial_zT_2&=\cos z/(18\sqrt Q),\\
 \partial_xT_3&=ye^{-xy}/20,&\partial_yT_3&=xe^{-xy}/20.
\end{align*}
Las de $T_1$ están acotadas por $\sin1/3<0.3$; las de $T_2$ por $1/(9\sqrt{0.21})<0.3$ y $1/(18\sqrt{0.21})<0.3$; las de $T_3$ por $e/20<0.3$. Cada fila tiene suma menor que $0.9$, de modo que $q=0.9$ es una constante válida. Banach garantiza un único punto fijo de esta aplicación en el cubo.

En este caso el punto fijo puede darse exactamente:
\[
 (x_*,y_*,z_*)=(1/2,0,-\pi/6).
\]
La primera coordenada de $T$ es $1/3+1/6=1/2$; para la segunda el radicando vale $1/4-1/2+1.06=0.81$, cuya raíz es $0.9$; la tercera es $-1/20-(10\pi-3)/60=-\pi/6$. Se comprueban así las tres igualdades. La unicidad afirmada corresponde a la aplicación y a la rama elegida; despejar la raíz positiva no permite descartar sin examen las soluciones de la otra rama del sistema original.
\end{ejemplo}

\section{La distancia de Hausdorff y un fractal como punto fijo}
Sea $\mathcal F(X)$ la familia de subconjuntos \emph{no vacíos}, cerrados y acotados de $X$. Para $A,B\in\mathcal F(X)$ definimos
\[
 H(A,B)=\max\left\{\sup_{a\in A}\dist(a,B),
                         \sup_{b\in B}\dist(b,A)\right\}.
\]
\index{Hausdorff, distancia de}\index{fractal}\index{Cantor, conjunto de}
La no vacuidad evita distancias indefinidas y supremos sobre familias vacías. La acotación hace finitos los dos supremos: fijando $b_0\in B$, por ejemplo, $\dist(a,B)\leq d(a,b_0)$ queda acotada sobre $A$.

\begin{proposicion}\label{ch11:hausdorff-metrica}
La función $H$ es una métrica en $\mathcal F(X)$. Además,
\[
 |\dist(x,A)-\dist(x,B)|\leq H(A,B)\qquad(x\in X).
\]
\end{proposicion}
\begin{proof}
La simetría y la no negatividad están en la definición. Si $H(A,B)=0$, cada $a\in A$ tiene distancia cero a $B$ y, como $B$ es cerrado, pertenece a $B$. Al intercambiar los conjuntos obtenemos $A=B$. La vuelta es inmediata.

Para la desigualdad triangular observamos que, dados $a\in A$ y $\eta>0$, podemos elegir $b\in B$ con $d(a,b)<\dist(a,B)+\eta$ y después $c\in C$ con $d(b,c)<\dist(b,C)+\eta$. Entonces
\[
 \dist(a,C)\leq\dist(a,B)+H(B,C)+2\eta.
\]
Dejando $\eta\downarrow0$ y tomando el supremo en $a$ queda una cota por $H(A,B)+H(B,C)$. La otra distancia dirigida, de $C$ hacia $A$, se acota escogiendo primero un punto de $B$ próximo a un punto de $C$ y después uno de $A$ próximo a ese punto de $B$; da la misma suma. Tomar el máximo prueba $H(A,C)\leq H(A,B)+H(B,C)$.

Para la última afirmación, elegimos $a\in A$ con $d(x,a)<\dist(x,A)+\eta$ y un punto de $B$ a distancia menor que $H(A,B)+\eta$ de $a$. Resulta $\dist(x,B)\leq\dist(x,A)+H(A,B)+2\eta$. Al dejar $\eta\downarrow0$ e intercambiar $A,B$ obtenemos la estimación.
\end{proof}

El ejercicio 5 demostrará que $\mathcal F(X)$ es completo exactamente cuando $X$ lo es. Para $X=\R$, definamos
\[
 \mathcal T(A)=\{a/3:a\in A\}\cup\{(a+2)/3:a\in A\}.
\]
Las transformaciones afines son homeomorfismos de $\R$, por lo que sus imágenes de cerrados son cerradas; también preservan no vacuidad y acotación. Así $\mathcal T$ actúa en $\mathcal F(\R)$. Para un punto $a/3$ de la primera pieza y un punto $b\in B$ próximo a $a$,
\[
 \dist(a/3,\mathcal T(B))\leq|a-b|/3.
\]
Tomando el ínfimo sobre $b$ se obtiene una cota $\dist(a,B)/3$. Para $(a+2)/3$ se usa la segunda pieza de $\mathcal T(B)$ y resulta la misma cota. Tomando supremos, y después intercambiando $A,B$, concluimos
\[
 H(\mathcal T(A),\mathcal T(B))\leq\tfrac13H(A,B).
\]
Por Banach existe un único conjunto fijo. Es el conjunto de Cantor: en su representación ternaria con dígitos $0$ y $2$, la primera pieza antepone el dígito $0$ y la segunda antepone el dígito $2$. Toda sucesión de esos dígitos comienza con uno de ellos, por lo que $C=C/3\cup(C+2)/3$. Las iteraciones desde cualquier cerrado acotado no vacío convergen a $C$ en distancia de Hausdorff.

\section{Ejercicios resueltos}
\begin{ejercicio}[Lectura 11, ejercicio 1]
Si $X$ es completo no vacío y $f:X\to X$ satisface que $f\circ f$ es una contracción, determine si $f$ tiene un punto fijo.
\end{ejercicio}
\begin{solucion}
La aplicación $T=f^2$ tiene un único punto fijo $a$. No necesitamos suponer que $f$ sea continua. Por asociatividad de la composición,
\[
 T(f(a))=f(f(f(a)))=f(T(a))=f(a).
\]
Así $f(a)$ también es punto fijo de $T$; por unicidad, $f(a)=a$. Si $b$ fuera otro punto fijo de $f$, sería punto fijo de $T$, luego $b=a$. En general, el mismo razonamiento funciona si alguna potencia $f^k$ es contracción: $f$ conmuta con sus potencias y todo punto fijo de $f$ es fijo para $f^k$.
\end{solucion}

\begin{ejercicio}[Lectura 11, ejercicio 2]
Usando el ejercicio anterior, pruebe que $x=\cos x$ tiene una única solución en $[0,1]$.
\end{ejercicio}
\begin{solucion}
La aplicación $f(x)=\cos x$ lleva $[0,1]$ a $[\cos1,1]\subset[0,1]$. Su cuadrado por composición es $f^2(x)=\cos(\cos x)$ y
\[
 |(f^2)'(x)|=|\sin(\cos x)\sin x|\leq\sin1<1.
\]
El valor medio prueba que $f^2$ es contracción en el intervalo completo; el ejercicio 1 da un único punto fijo de $f$. También se podría aplicar Banach directamente porque $|f'(x)|=\sin x\leq\sin1<1$. La formulación mediante $f^2$ ilustra que la propiedad contractiva puede verificarse tras varias iteraciones.
\end{solucion}

\begin{ejercicio}[Lectura 11, ejercicio 3]
Pruebe que $f(x)=x^2+x^{-1}$ no tiene puntos fijos en $[1,\infty)$.
\end{ejercicio}
\begin{solucion}
Para $x\geq1$, $x^2\geq x$ y $x^{-1}>0$, de donde $f(x)>x$. Por tanto $f(x)=x$ es imposible. La aplicación sí lleva el intervalo a sí mismo, pero no es una contracción: su derivada $2x-x^{-2}$ es no acotada. La invariancia de un espacio completo, sin contracción, no garantiza un punto fijo.
\end{solucion}

\begin{ejercicio}[Lectura 11, ejercicio 4]
Sean $0<r<1$ y $\alpha\in\R$. En $C([0,r])$ defina
\[
 (\Phi f)(x)=\alpha+\int_0^x f(t)\,dt.
\]
Pruebe que $\Phi$ es contracción para la norma uniforme y encuentre su punto fijo.
\end{ejercicio}
\begin{solucion}
La integral de una función continua tiene una primitiva continua, por lo que $\Phi$ toma valores en el espacio indicado. Además,
\[
 |(\Phi f)(x)-(\Phi g)(x)|
 \leq\int_0^x|f(t)-g(t)|\,dt
 \leq r\norm{f-g}_\infty.
\]
Tomando el supremo resulta $\norm{\Phi f-\Phi g}_\infty\leq r\norm{f-g}_\infty$, con $r<1$.

Justificamos la completitud del dominio sin anticipar resultados: una sucesión de Cauchy en norma uniforme tiene un límite puntual porque $\R$ es completo. Pasando al límite en una de las dos funciones de la estimación uniforme de Cauchy, se obtiene convergencia uniforme al límite. Ese límite es continuo: dado $x$ y $\eps>0$, elegimos una función de la sucesión a distancia uniforme menor que $\eps/3$ del límite, y usamos su continuidad en $x$ con tolerancia $\eps/3$; la desigualdad triangular suma tres errores menores que $\eps/3$. La función límite es acotada por una función de la sucesión más un error uniforme finito. Esto prueba que $C([0,r])$ es completo.

Banach proporciona un único punto fijo. La función $f(x)=\alpha e^x$ lo es, pues
\[
 \alpha+\int_0^x\alpha e^t\,dt=\alpha+\alpha(e^x-1)=\alpha e^x.
\]
También se descubre la solución diferenciando la ecuación integral: un punto fijo satisface $f'=f$ y $f(0)=\alpha$; multiplicar por $e^{-x}$ muestra que la derivada de $e^{-x}f(x)$ es cero.
\end{solucion}

\begin{ejercicio}[Lectura 11, ejercicio 5]\label{ch11:hiper-completo}
Demuestre que $(\mathcal F(X),H)$ es completo si y sólo si $X$ es completo.
\end{ejercicio}
\begin{solucion}
Supondremos primero $X$ completo y no vacío. Sea $(A_n)$ de Cauchy para $H$. Las funciones $u_n(x)=\dist(x,A_n)$ satisfacen
\[
 |u_n(x)-u_m(x)|\leq H(A_n,A_m)
\]
para todo $x$. Por completitud de $\R$, existe $u(x)=\lim_nu_n(x)$, y la estimación de Cauchy, al pasar $m\to\infty$, demuestra que la convergencia es uniforme en todo $X$. No es necesario que las funciones $u_n$ sean acotadas individualmente. Cada una es $1$-Lipschitz; al tomar límites en $|u_n(x)-u_n(y)|\leq d(x,y)$ vemos que $u$ también lo es. Por tanto
\[
 A=\{x\in X:u(x)=0\}
\]
es cerrado. Demostraremos que es no vacío y acotado, y que $H(A_n,A)\to0$.

\textbf{Construcción de un punto límite.} Seleccionemos índices crecientes $n_k$ con $H(A_{n_k},A_{n_{k+1}})<2^{-k-1}$. Elegimos $a_1\in A_{n_1}$ y sucesivamente $a_{k+1}\in A_{n_{k+1}}$ con $d(a_k,a_{k+1})<2^{-k}$. Tal elección es posible por la definición de distancia a un conjunto y la cota estricta anterior. Como $\sum_k2^{-k}<\infty$, las colas de la sucesión satisfacen $d(a_l,a_k)\leq\sum_{j=k}^{l-1}2^{-j}$; es de Cauchy y converge a algún $a\in X$. La convergencia uniforme $u_{n_k}\to u$ y $u_{n_k}(a_k)=0$ implican
\[
 |u(a)|\leq |u(a)-u(a_k)|+|u(a_k)-u_{n_k}(a_k)|\longrightarrow0.
\]
Así $a\in A$, y $A$ no es vacío.

\textbf{Acotación.} Para algún índice $N$ y alguna constante finita $M$, la convergencia uniforme da $|u_N(x)-u(x)|\leq M$ para todo $x$. Si $x\in A$, entonces $\dist(x,A_N)\leq M$. Como $A_N$ es acotado, elegimos $a_0\in A_N$ y $R$ con $d(a_0,b)\leq R$ para todo $b\in A_N$. Para cualquier $x\in A$ existe $b\in A_N$ con $d(x,b)<M+1$, así que $d(x,a_0)<M+1+R$. Esto prueba que $A\in\mathcal F(X)$.

\textbf{Convergencia en ambas direcciones.} Dado $\eps>0$, elegimos $N$ tal que $H(A_n,A_m)<\eps$ para $n,m\geq N$. Al dejar $m\to\infty$ en la desigualdad para las funciones de distancia, obtenemos $|u_n(x)-u(x)|\leq\eps$ para todo $x$ y $n\geq N$. En particular,
\[
 \sup_{a\in A}\dist(a,A_n)\leq\eps.
\]
Para la dirección restante fijemos $n\geq N$ y $b\in A_n$. Construiremos una cadena que empieza en $b$. Para cada $k\geq0$ elijamos umbrales crecientes $N_k$, con $N_0=N$, tales que $H(A_l,A_m)<\eps2^{-k}$ para $l,m\geq N_k$. Pongamos $n_0=n$, y para $k\geq1$ escojamos índices crecientes $n_k\geq N_k$. Entonces
\[
 H(A_{n_k},A_{n_{k+1}})<\eps2^{-k}\quad(k\geq0).
\]
Partiendo de $b_0=b$, elegimos $b_{k+1}\in A_{n_{k+1}}$ con $d(b_k,b_{k+1})<2\eps2^{-k}$. La suma de estas cotas es $4\eps$, de modo que $(b_k)$ es de Cauchy y converge a un punto $b_\infty$ con $d(b,b_\infty)\leq4\eps$. El mismo argumento con $u_{n_k}$ usado antes demuestra $b_\infty\in A$. Así $\dist(b,A)\leq4\eps$, uniformemente para $b\in A_n$. Concluimos $H(A_n,A)\leq4\eps$ para $n\geq N$. Como $\eps$ es arbitrario, hay convergencia.

Recíprocamente, supongamos completo $\mathcal F(X)$ y sea $(x_n)$ de Cauchy en $X$. La identidad $H(\{x_n\},\{x_m\})=d(x_n,x_m)$ muestra que los singletons convergen en $\mathcal F(X)$ a algún $A\ne\varnothing$. Para $a,b\in A$,
\[
 d(a,b)\leq d(a,x_n)+d(x_n,b)\leq2H(A,\{x_n\})\longrightarrow0.
\]
Así $A$ tiene un solo punto, digamos $x$, y $d(x_n,x)=H(\{x_n\},A)\to0$. Por tanto $X$ es completo. Si $X=\varnothing$, también $\mathcal F(X)=\varnothing$ y ambos espacios son completos por la definición; la equivalencia sigue siendo válida.
\end{solucion}

\begin{ejercicio}[Lectura 11, ejercicio 6]
Calcule por iteración la solución del sistema
\[
 5x^2-y^2=0,\qquad y-\tfrac14(\sin x+\cos y)=0
\]
que se encuentra en la región positiva próxima a $(1/4,1/4)$.
\end{ejercicio}
\begin{solucion}
Para la rama positiva $y=\sqrt5\,x$, tomamos
\[
 T(x,y)=\left(\frac{y}{\sqrt5},\frac{\sin x+\cos y}{4}\right)
\]
en el rectángulo completo $R=[0,1/2]^2$. La primera coordenada está entre cero y $1/(2\sqrt5)<1/2$. La segunda está entre $\cos(1/2)/4>0$ y $(\sin(1/2)+1)/4<1/2$. Por tanto $T(R)\subset R$. Las sumas de valores absolutos de sus derivadas por fila son $1/\sqrt5<1/2$ y $(|\cos x|+|\sin y|)/4\leq(1+\sin(1/2))/4<1/2$. Podemos usar $q=1/2$ como constante de contracción.

La iteración simultánea, empezando en $(x_0,y_0)=(0.25,0.25)$, es
\[
 x_{n+1}=y_n/\sqrt5,\qquad y_{n+1}=(\sin x_n+\cos y_n)/4.
\]
En ambas fórmulas se usan los valores de la iteración anterior. Algunos valores calculados son:
\begin{center}
\begin{tabular}{rrr}
\toprule
$n$&$x_n$&$y_n$\\\midrule
0&0.250000000000&0.250000000000\\
1&0.111803398875&0.304079095241\\
2&0.135988305499&0.266423427535\\
3&0.119148178954&0.275072068197\\
5&0.120889895207&0.271597989694\\
10&0.121245551908&0.271102661689\\
15&0.121241889930&0.271105174792\\
20&0.121241911634&0.271105155669\\\bottomrule
\end{tabular}
\end{center}
La cota a priori da
\[
 \norm{(x_n,y_n)-(x_*,y_*)}_\infty
 \leq 2^{1-n}\norm{(x_1,y_1)-(x_0,y_0)}_\infty
 <0.277\,2^{-n}.
\]
En la iteración $20$, el error garantizado es menor que $2.65\cdot10^{-7}$, sin confundir los decimales mostrados con todos los decimales certificados. Así la solución es aproximadamente $(0.121242,0.271105)$, única dentro del rectángulo para la rama elegida. Sustituir el punto fijo en las dos fórmulas demuestra que resuelve el sistema original.
\end{solucion}

\section{Tres comprobaciones antes de iterar}
La aplicación debe conservar el dominio, ese dominio debe ser completo y la constante de contracción debe ser estrictamente menor que uno. Verificados esos tres puntos, la misma demostración produce soluciones de ecuaciones, funciones definidas por ecuaciones integrales y conjuntos geométricos autosimilares, con cotas de error explícitas.
